\documentclass[12pt]{article}
\usepackage[T1]{fontenc}
\usepackage[utf8]{inputenc}
\usepackage{lmodern}
\usepackage{textcomp}
\usepackage{enumitem}
\usepackage{geometry}
\geometry{margin=1in}
\begin{document}
\begin{center}
Answer Key - Computer Science - K-12 Level Assessment 24 \
\small (Answers are ordered exactly as in the question paper.)
\end{center}

\section*{Multiple Choice}
\begin{enumerate}[label=\arabic*.]
\item \textbf{Answer:} B
\\ \textit{Options:} A) replace(old, new [, max]); B) strip([chars]); C) swapcase(); D) title()
\\ \textit{Note:} The `strip([chars])` function in Python 3 is designed specifically to remove all leading and trailing whitespace from a string, making it the correct answer for this question.
\item \textbf{Answer:} A
\\ \textit{Options:} A) 1; B) 2; C) 3; D) 4
\\ \textit{Note:} The correct answer is 1 because the `min()` function in Python 3 returns the smallest element in a list, and in the list l = [1, 2, 3, 4], the smallest value is 1.
\item \textbf{Answer:} A
\\ \textit{Options:} A) O(1); B) O(n); C) O($n^2$); D) O(log n)
\\ \textit{Note:} O(1) is considered the smallest asymptotic notation because it represents a constant time complexity that does not grow with the size of the input, making it the most efficient in terms of performance.
\item \textbf{Answer:} B
\\ \textit{Options:} A) A device driver is assigned to the device.; B) An Internet Protocol (IP) address is assigned to the device.; C) A packet number is assigned to the device.; D) A Web site is assigned to the device.
\\ \textit{Note:} When a new device connects to the Internet, it is assigned an Internet Protocol (IP) address, which uniquely identifies it on the network and enables communication with other devices.
\item \textbf{Answer:} C
\\ \textit{Options:} A) O(1); B) O(n); C) O($n^2$); D) O(log n)
\\ \textit{Note:} O($n^2$) is the largest asymptotic notation among the options provided because it grows faster than linear (O(n)) or logarithmic (O(log n)) functions as the input size n increases, indicating that it represents a higher rate of growth in computational complexity.
\end{enumerate}

\section*{True / False}
\begin{enumerate}[label=\arabic*.]
\setcounter{enumi}{5}
\item \textbf{Answer:} True
\\ \textit{Options:} A) True; B) False
\\ \textit{Note:} The sum of the list l = [1, 2, 3, 4] in Python 3 is 10 because adding the numbers together (1 + 2 + 3 + 4) equals 10.
\item \textbf{Answer:} True
\\ \textit{Options:} A) True; B) False
\\ \textit{Note:} In Python 3, the expression x \textless{}\textless{} 3 performs a left bitwise shift on the binary representation of x, which multiplies it by 2 raised to the power of 3; therefore, when x is 1 (binary 0001), shifting it left by 3 results in 8 (binary 1000).
\item \textbf{Answer:} False
\\ \textit{Options:} A) True; B) False
\\ \textit{Note:} The statement is false because the expression ['Hi!'] * 4 successfully creates a list containing four references to the string "Hi!" without raising any syntax errors.
\item \textbf{Answer:} False
\\ \textit{Options:} A) True; B) False
\\ \textit{Note:} Averaging numeric values may mask outliers and inaccuracies, making it an ineffective method for detecting data entry errors, as it does not identify individual discrepancies in the dataset.
\item \textbf{Answer:} False
\\ \textit{Options:} A) True; B) False
\\ \textit{Note:} The statement is false because O(n) represents a linear growth rate, which increases more slowly than O($n^2$), a quadratic growth rate, as the input size increases.
\end{enumerate}

\section*{Long Form Response}
\begin{enumerate}[label=\arabic*.]
\setcounter{enumi}{10}
\item \textbf{Answer:} def series\_sum(number): | return total
\\ \textit{Note:} The answer correctly identifies the need for a function that calculates the sum of squares of the first n integers, although it lacks the implementation details to compute the total.
\item \textbf{Answer:} def chinese\_zodiac(year): | return sign
\\ \textit{Note:} correct because it outlines the structure of a function named `chinese\_zodiac` that takes a year as input and is intended to return the corresponding zodiac sign, capturing the essential requirement of the question.
\end{enumerate}

\vfill
\noindent\textit{End of Answer Key}
\end{document}